\chapter[Band 24: Schlichte ungerichtete Graphen]{Band 24 auf einen Blick}
\noindent\textbf{Schlichte ungerichtete Graphen}\par\medskip
Ausgangspunkt sind ungerichtete relationale Graphen.
Schlichtheit ergänzt ausschließlich die Irreflexivität; Zusammenhang
gehört zur nächsten eigenständigen Strukturstufe.

\begin{BandUebersicht}
\textbf{Schlichtheit als Zusatzaxiom}
\UebFormel{\begin{aligned}
&\SimpleUndirectedGraphStruct{V,E}\ \leftrightarrow\\
&\quad\UndirectedGraphStruct{V,E}\land\\
&\quad\forall x\in V\quad(x,x)\notin E.
\end{aligned}}
\UebQuellen{\FormulaRefAuto{SimpleUndirectedGraphDef}[def]}
&
\textbf{Keine Knotenschleifen}
\UebFormel{\begin{aligned}
&\SimpleUndirectedGraphStruct{V,E},\
 x\in V\ \vdash\\
&\qquad(x,x)\notin E.
\end{aligned}}
\UebQuellen{\FormulaRefAuto{SimpleUndirectedGraphLoopFree}[thm]} \\
\addlinespace[.8ex]

\textbf{Einschränkung auf eine Knotenmenge}
\UebFormel{\begin{aligned}
&U\subseteq V,\qquad E[U]=E\cap(U\times U).
\end{aligned}}
\UebQuellen{\FormulaRefAuto{DirectedInducedEdgeRelationDef}[def]}
&
\textbf{Induzierte Teilgraphen sind wieder schlicht}
\UebFormel{\begin{aligned}
&\SimpleUndirectedGraphStruct{V,E},\
 U\subseteq V\ \vdash\\
&\qquad\SimpleUndirectedGraphStruct{U,E[U]}.
\end{aligned}}
\UebQuellen{\FormulaRefAuto{SimpleUndirectedInducedSubgraph}[thm]} \\
\addlinespace[.8ex]
\end{BandUebersicht}

